\documentclass[script,pro]{debate}
\motion{AI的迅猛发展提升 / 降低了人类创作者存在的意义}
\stance{正方}{提升了}
\meta{赛制}{招新赛 3v3}
\meta{口径来源}{备赛文档 第十节 正方口径表}
\begin{document}
\debatetitle
\begin{thesisbox}[全场一句话立论]机器能做的越多，“这是一个人做的”就越被区分、被看重，作品的分量也越落在人的立意上。\end{thesisbox}

\begin{speech}{1. 正一开篇立论稿}{正文 825 字 / 180 秒}
谢谢主席，各位好。先说清三个词。

AI 的迅猛发展，指的是 2022 年以来，机器生成文字、图像、音乐和声音的能力突然变强，而且人人可用。人类创作者，是把创造文艺作品当成自己一部分的人，靠它吃饭的算，业余在写在画的也算。存在的意义，词典里叫作用与重要性。放到今天是两个问题：这个世界还有多需要、多珍视由人来创作；创作还能给创作者自己多少东西。

所以判准很简单。和 AI 爆发之前比，由人来创作这件事，整体上变重了还是变轻了。被需要和被珍视，两半都要算；而且要看多数创作者，不能只看最红的几个人，也不能只看最惨的几个人。我方要证明它变重了，对方要证明它变轻了。只证明了“有”，都不算完成任务。

我方第一点，人味成了稀缺品。你想想，一张图几秒钟就能生成，要多少有多少，那稀缺的就不再是作品，而是作品背后那个有经历、要负责的人。心理学家做过一个实验：同一批 AI 画的画，随机贴上“人类创作”或者“AI 创作”的标签。贴了人类标签的，喜欢、美感、深刻、价值，四项评分全都更高。这份看重，社会正在把它写成规则。2025 年，美国上诉法院判定版权法里的作者必须是人，2026 年最高法院拒绝再审。我国从 2025 年 9 月起，要求 AI 生成的内容必须标识。皮尤研究中心的调查里，七成六的美国成年人认为，能分辨作品出自人还是 AI 很重要。\stress{机器越多，社会越要把人认出来。}

我方第二点，立意回到人。AI 接走的是画得像、写得顺，留给人的是想画什么、选哪一张。北京互联网法院 2023 年判过一个案子：有人用 AI 生成图片，法院认定，他设计提示词、设置参数，对画面做了个性化的选择和安排，体现了独创性的智力投入，著作权归他。说白了，法律认的是人的选择。一项追踪五万多名用户、四百万件作品的研究发现，用上 AI 以后，产量多了四分之一，每次被浏览时获赞的概率高了一半；得分更高的，是敢于探索新想法的人。门槛降下来，有话要说的人也更多了：另一项写作实验里，AI 的点子让原本最不会写故事的人进步最大。

我方承认，一部分执行型的订单确实在减少。但今天的题目问的是意义，不是订单数。所以我们全场都会请对方回答一个问题：如果人的来源不重要，为什么法院、行业和七成以上的普通人，都坚持要把人和机器分开？

机器能做的越多，“这是一个人做的”就越重。谢谢。
\end{speech}
\begin{contingency}
\ifthen{反方把意义定义成收入与岗位}{意义是作用与重要性，被需要和被珍视两半都算；照反方的定义，一个没有稿费的诗人没有存在的意义，这不是常识。}
\ifthen{反方说“标签只是防骗”}{想知道背后是不是一个人，本身就是在乎来源；人不会要求给流水线零件标上工人的名字。}
\end{contingency}

\begin{stage}{2. 正一被质询防守预案}{反二质询，90 秒双边计时}
\textbf{原则}：先给是或否，再用一句话守住前提；不否认常识，不拖时间（故意不答会扣分）。反二是全场第一次质询，最想拿到的是“执行型订单在减少”这句承认：可以承认，但一定接上“题目问的是意义，不是订单数”。
\begin{qa}
\qarow{插画师因为 AI 丢了活，贵方承认吗？}{承认，执行型的订单在减少。但丢一份活，不等于人来创作这件事变轻。}
\qarow{被需要算不算存在意义的一部分？}{算，是一半。另一半是被珍视，判准说两半一起算。}
\qarow{那被需要这一半，是变重还是变轻？}{“给我一张图”变轻了；“给我一个人做的图”变重了，法院和行业都在单独保护它。}
\qarow{给作品标上“人做的”，能带来订单吗？}{被认出来，是被选择的前提。我方说的是珍视，不是价格。}
\qarow{被珍视却没人来请，算意义提升吗？}{算。意义是作用与重要性，不是接单量；一个没有稿费的诗人，也有存在的意义。}
\qarow{那个贴标签的实验，画全是 AI 画的吧？}{对。画一模一样，评分的差别只能来自“是不是人做的”，这正是实验要测的。}
\qarow{获赞那项研究，样本是 AI 绘画网站的用户吧？}{是，这个局限我方承认。我们用它说明分数给了立意，不用它代表所有画师。}
\end{qa}
\end{stage}

\begin{stage}{3. 正二质询反一问题链}{90 秒，双边计时}
\textbf{总目标}：让反一承认“被珍视”也算意义，并承认决定作品归属的是人的选择。两条链共五问，对方答非所问先复述一次问题，第二次才礼貌打断。
\begin{chain}{链一：两半链}{让对方承认意义包括被珍视}
\q{贵方说的意义，包不包括被珍视？}{答“包括”：下一问；答“不包括”：追“那贵方的意义就只剩订单，对吗？”}
\q{法院只认人当作者，算不算看重人？}{答“算”：下一问；答“不算”：追“被法律单独保护，也不算被看重吗？”}
\q{珍视在变重，贵方还说整体变轻，对吗？}{答“对”：收束，点出对方只算了一半；答“珍视落在少数”：转链二}
\cut{好，我换个问法。}
\close{感谢，对方承认被珍视也是意义的一部分。}
\end{chain}
\begin{chain}{链二：选择链}{让对方承认决定作品归属的是人的选择}
\q{北京那个 AI 图的案子，判给人还是 AI？}{答“人”：下一问}
\q{法院看的，是不是人的选择和安排？}{答“是”：收束；答“那是提示词”：追“提示词是不是人写的？”}
\close{感谢。法院认的是人的选择，这就是我方第二点。}
\end{chain}
\end{stage}

\begin{stage}{4. 正三对辩要点}{双方各 90 秒}
\textbf{任务}：开第一个战场：意义是两半，需要和珍视都算。每次发言先一句回应，再抛一问。
\begin{points}{进攻点（4–5 个，每个 15–20 秒）}
\item 判准两半：被需要和被珍视都算。贵方的数据全是订单，珍视那一半算过没有？
\item 法院只认人当作者，我国要求 AI 内容标识：这是不是社会在区分并看重人？
\item 同一批 AI 画，贴上人类标签四项评分全高：人们看重的到底是什么？
\item 写作类月订单降了百分之二，这能叫由人来创作整体变轻吗？
\item 七成六的美国成年人认为能分辨人或 AI 很重要：想分辨，是不是在乎？
\end{points}
\begin{points}{备用点（6–8 个）}
\item 对方说被区分不等于被需要 → 区分是被选择的前提，认不出的东西谁也选不了。
\item 对方说珍视只流向少数 → 标识约束的是机器，不挑名气；法院认的是选择，不是粉丝数。
\item 对方说标识办法是为了防虚假信息 → 目的是什么都好，结果是“是不是人做的”成了必须公开的事实。
\item 对方说九成七的人分不出来 → 同一份调查里一半以上的人为此不安，八成要求标注。
\item 对方说保护说明濒危 → 意义是重要性，不是安全感；人只为自己看重的东西立法。
\item 对方说顶尖者受冲击更大 → 那说明量到的是执行的订单，不是人的分量。
\item 对方说到 2028 年收入面临风险 → 那是预测，而且是收版税的机构做的。
\end{points}
\begin{points}{对方最可能问的三个问题 → 回应}
\item 标签能带来订单吗？ → 我方说珍视不说价格；被认出来是被选择的前提。
\item 丢了活的插画师，意义提升了吗？ → 她的处境变差了，我方承认；但人来创作这件事变重，不等于每一个人都更好过。
\item 分不出来的珍视有什么用？ → 所以才要标注。八成的人要求标注，说明大家想把人认出来。
\end{points}
\end{stage}

\begin{kit}{5. 正二申论}{套件 · 组装后 396–416 字 / 90 秒}
\assembly{开头 → 回应反二（任选 1） → 拆反方立论（任选 1） → 推进论点一 → 收尾}
\begin{fixed}{开头}
谢谢主席，各位好。我做两件事：先看反方立论站不站得住，再把我方第一点往前推一步。
\end{fixed}
\begin{pick}{回应反二}{1}
\begin{module}{如果反二申论主打「被区分不等于被需要」}
反二刚才说，被区分不等于被需要。这两个词确实不一样。但题目问的是意义，意义本来就包括被珍视。如果把珍视这一半拿掉，只剩接单，那就不是在比意义了。
\end{module}
\begin{module}{如果反二申论主打「立意回到的是甲方」}
反二刚才说，立意回到的是甲方。那我问一句，甲方是不是人？出于己意去创造的人，就是创作者。立意从画师手里到了甲方手里，也还在人手里，没有到机器手里。
\end{module}
\begin{module}{如果反二申论主打「订单和收入在下降」}
反二刚才把重点放在订单和收入上。我方承认，执行型的订单在减少。但订单和收入能证明的，是有人丢了活；还证明不了人来创作这件事变轻了，这是两回事。
\end{module}
\begin{fallback}{反二申论的重点不在以上几条时}
先回到判准。今天比的是由人来创作这件事变重还是变轻，需要和珍视两半都要算，还要看多数创作者，不看最红的，也不看最惨的。我方就按这个标准往下说。
\end{fallback}
\end{pick}
\begin{pick}{拆反方立论}{1}
\begin{module}{如果反方立论建立在「订单和收入的下降」上}
反方立论的主干，是订单和收入在下降。我方看三处。第一，量级：最常被引的那项平台研究里，写作类的月订单降了百分之二，能叫整体变轻吗？第二，性质：常见的另外两份，一份是会员自愿填的问卷，一份是收版税的机构做的预测，都不是测量。第三，最要紧：订单只是按件计价的那一半。社会一边少买执行，一边在立法、在标注、在评价里把人的作品单独挑出来，这一半同样要算。
\end{module}
\begin{module}{如果反方立论建立在「作品里属于人的部分变少」上}
反方立论的主干，是作品里属于人的部分在变少。我方看两处。第一，一个插画师被请去修图、被压价，是市场给执行开的价变低了，不是她的立意不值钱了；想画什么，机器替不了任何人。第二，最常被引的那项五万多名用户的研究，同一篇里还写着：敢于探索新想法的人，被评价得更高。分数给了立意，没有给平均。只念前半句，就漏掉了后半句。
\end{module}
\begin{module}{如果反方立论建立在「观众分不出人和机器」上}
反方立论的主干，是观众分不出来。我方看两处。第一，分不出来，说明机器学得像，不说明人不重要。最常被引的那份八国九千人的盲听调查里，九成七的人分不出来；可同一份调查里，超过一半的人为此感到不安，八成要求给纯 AI 音乐清楚标注。第二，认不出是一个信息问题，这个问题正在被标识制度解决。人会为分不出来而不安，恰恰说明大家在乎背后是不是一个人。
\end{module}
\begin{fallback}{反方立论的主干不在以上几条时}
不管反方立论用什么论据，都绕不开一个问题：它证明了有人受损，还没有证明整体变轻。判准要两半一起算，还要看多数。反方要赢，得证明多数创作者被需要和被珍视加在一起变少了。按件计价的订单和收入，只是其中一半，另一半是法律、行业和观众给人的看重。只拿出其中一半，这笔账就还没有算完，结论也就还下不了。
\end{fallback}
\end{pick}
\begin{fixed}{推进论点一}
我方第一点再往前走一步。人味成了稀缺品，不只写在法律里，也写进了行业的规矩。2023 年好莱坞编剧罢工之后，新合同写明：AI 不能写、也不能改写剧本，公司不能要求编剧用 AI。2025 年，美国作家协会推出“人类创作”认证，给人写的书打上可查的标记。说白了，机器能写以后，“这是人写的”成了值得被证明的事。
\end{fixed}
\begin{fixed}{收尾}
所以我问一句：如果人的来源不重要，编剧、作家和法院为什么都在抢着把它写清楚？谢谢。
\end{fixed}
\end{kit}

\begin{stage}{6. 正二被质询防守预案}{反三质询，90 秒双边计时}
\textbf{原则}：正二刚推进了行业规矩这一层，反三最可能追“防骗不等于珍视”和“规矩只保护少数”。先答是或否，再一句话守住：我方说的是珍视，不是价格。
\begin{qa}
\qarow{标识办法写的目的，是防虚假信息吧？}{是。目的是什么都好，结果是“是不是人做的”成了必须公开的事实。}
\qarow{防骗的标签，也能算对画师的珍视？}{想知道背后是不是一个人，本身就是在乎来源。没人要求给流水线零件标上工人的名字。}
\qarow{九成七的人盲听分不出来，对吗？}{对。分不出来说明机器学得像；同一份调查里八成要求标注，说明大家想把人认出来。}
\qarow{编剧合同只保护工会里的编剧吧？}{合同是签给工会会员的。但它第一次把“AI 不能当编剧”写成行业规矩，说明行业认定编剧这件事得是人来做。}
\qarow{被珍视却没人来请，这算提升？}{算。意义是作用与重要性，不是接单量；而且被认出来，是被选择的前提。}
\qarow{立意是不是回到了甲方？}{甲方也是人。立意在谁手里都在人手里，没有到机器手里。}
\end{qa}
\end{stage}

\begin{stage}{7. 正三质询反二问题链}{90 秒，双边计时}
\textbf{总目标}：压反二数据的量级，并让反二承认“要求标注是在乎来源”。四问，不恋战。
\begin{chain}{链一：量级链}{让对方承认订单下降的幅度撑不起“整体变轻”}
\q{贵方引的研究，写作类月订单降百分之二，对吗？}{答“对”：下一问；答“还有收入”：追“收入降百分之五以上，对吗？”再下一问}
\q{这个幅度，能证明整体变轻吗？}{答“能”：收束，请评委自己判断；答“还有别的数据”：转链二}
\cut{好，我换个问法。}
\close{感谢，对方的主干数据是个位数的下降。}
\end{chain}
\begin{chain}{链二：珍视链}{让对方承认要求标注是在乎来源}
\q{八成受访者要求给纯 AI 音乐标注，贵方否认吗？}{答“不否认”：下一问}
\q{要求标注，是不是在乎背后是不是人？}{答“是”：收束；答“是怕被骗”：追“怕被骗，是不是在乎谁做的？”}
\close{感谢。在乎背后是不是人，就是我方说的珍视。}
\end{chain}
\end{stage}

\begin{stage}{8. 正一对辩要点}{双方各 90 秒}
\textbf{任务}：在中段收拢前两段的交锋，把评委拉回判准：两半都算，看多数。
\begin{points}{进攻点（4–5 个，每个 15–20 秒）}
\item 判准两半：贵方到现在为止，珍视那一半算过没有？
\item 看多数：标识和法律对每一个人类创作者都适用，不挑名气。这算不算多数？
\item 那项五万多名用户的研究里，获赞集中度下降了：好处是不是不只给了头部？
\item 执行变便宜了，立意是不是更值钱了？
\item 贵方的数据量的是订单，贵方的题目问的是意义，两者是一回事吗？
\end{points}
\begin{points}{备用点（6–8 个）}
\item 对方说那是少数 → 判准看多数，所以我方用的是法律和标注，它们对所有人一样。
\item 对方说修图工 → 被压价的是执行，想画什么机器替不了。
\item 对方说作品越来越像 → 越趋同，有立意的人越显眼。
\item 对方说珍视不当饭吃 → 题目问的是意义，不是饭碗；我方承认处境变难。
\item 对方说分不出来 → 分不出来的人里一半以上不安，八成要求标注。
\item 对方说顶尖者也受冲击 → 那说明量到的是执行的订单。
\end{points}
\begin{points}{对方最可能问的三个问题 → 回应}
\item 珍视落在谁身上？ → 落在每一件被认出来的人的作品上，标识和法律不挑名气。
\item 贵方承认订单在减少吗？ → 承认执行型订单在减少；题目问的是意义。
\item 甲方算创作者吗？ → 出于己意去创造的人都是创作者；但我方论点讲的是人在作品里的分量，不靠这一条。
\end{points}
\end{stage}

\begin{kit}{9. 正三申论}{套件 · 组装后 386–396 字 / 90 秒}
\assembly{开头 → 处理反方最强点（任选 1） → 推进论点二 → 收尾}
\begin{fixed}{开头}
谢谢主席，各位好。我做两件事：先处理反方最有力的一点，再把我方第二点往前推一层。
\end{fixed}
\begin{pick}{处理反方最强点}{1}
\begin{module}{如果反方全场最有力的是「修图工」}
反方最有力的一点，是修图工：人被压到只能给机器收尾。我方不回避，这很疼。但请评委分清两件事。被拿走的是按星期计价的执行，不是想画什么的能力。北京互联网法院判案时认的就是这一层：谁做了个性化的选择和安排，作品就归谁。出路不在跟机器比手快，而在把立意拿回来；今天一个人只要有立意，不用养一个团队也能把它做出来。
\end{module}
\begin{module}{如果反方全场最有力的是「订单和收入数据」}
反方最有力的一点，是订单和收入的数据。我方承认，执行型的订单在减少，写字的、画图的都有，这是真的。但请评委问一句：少掉的是哪一种需要？是“给我一张图”，不是“给我一个人做的图”。后一种需要，正在被法律、行业和标注单独保护起来。反方的数据只量到了前一种。判准要两半一起算，只量一半，结论就只能写“有人受损”，写不出“整体变轻”。
\end{module}
\begin{module}{如果反方全场最有力的是「分不出人和机器」}
反方最有力的一点，是大多数人分不出人和机器。我方承认，盲听盲读确实分不出来，八国九千人里九成七的人都分不出。但同一份调查里，超过一半的人为分不出来感到不安，八成要求标注。你想想，没有人会为分不清两颗螺丝钉而不安，人只会为分不清自己在乎的东西而不安。分不出来，说明机器学得像；要求分出来，才说明人的来源重要。
\end{module}
\begin{fallback}{反方最有力的一点不在以上几条时}
不管反方最有力的是哪一点，请评委用判准去量它。它说的是多数创作者，还是少数人？它量的是需要和珍视两半，还是只有一半？如果它证明的是有人受损，受损是真的，我方从不否认，那些丢了活的人值得被看见。可判准问的是整体方向，方向要两半一起算；只看需要那一半，算不出整体变轻，也算不出多数人的处境。
\end{fallback}
\end{pick}
\begin{fixed}{推进论点二}
我方第二点，立意回到人，再加一层。那项追踪五万多名用户的研究还发现，用了 AI 以后，获赞在用户之间的集中度下降了。也就是说，好处没有只流向头部，更多普通人的立意被看见了，这正是判准要看的多数。再给各位一个画面。美国乡村歌手 Randy Travis，2013 年中风，从此几乎说不出话。2024 年，团队用 AI 从他过去的四十二段录音里还原了他的声音，他本人和合作多年的制作人一起把关，发了一首新歌。他的妻子说，音乐就是他的全部。
\end{fixed}
\begin{fixed}{收尾}
机器给了他声音，歌还是他的。所以请反方回答：一个人重新能创作了，他存在的意义是变轻了吗？谢谢。
\end{fixed}
\end{kit}

\begin{stage}{10. 正三被质询防守预案}{反一质询，90 秒双边计时}
\textbf{原则}：反一会追“立意到底在谁手里”和“个例能不能代表多数”。个例坦白说是个例，普遍性交给五万多名用户的数据和法院规则。
\begin{qa}
\qarow{海报的立意，现在是甲方出还是画师出？}{谁出都是人出。出于己意去创造的人就是创作者；我方讲的是人在作品里的分量。}
\qarow{画师只负责修光线，算立意回到画师吗？}{那是执行被压价，我方承认。但她想画什么的能力机器拿不走，而且今天她自己的立意不用靠一个团队也能做出来。}
\qarow{同一研究里视觉新颖度全面下降，对吗？}{对，说的是平均。同一篇研究里，探索新想法的人被评价更高，分数给了立意。}
\qarow{Randy Travis 是明星个例吧？}{是个例，我方只用它说明机制。普遍性靠五万多名用户、四百万件作品的数据。}
\qarow{获赞集中度下降，也只是 AI 网站用户吧？}{是，这个局限我方承认。它说明在用 AI 的人里，好处没有集中到头部。}
\qarow{作品越来越像，属于作者的部分多了还是少了？}{平均越像，有立意的人越显眼；这正是作品的分量落回立意的另一面。}
\end{qa}
\end{stage}

\begin{stage}{11. 正一质询反三问题链}{90 秒，双边计时}
\textbf{总目标}：为结辩取材。拿到“决定画什么的是人”和“好处不只给头部”两句承认。
\begin{chain}{链一：立意链}{让对方承认立意始终在人手里}
\q{决定一张海报画什么的，是人还是 AI？}{答“人”：下一问；答“甲方”：追“甲方是不是人？”再下一问}
\q{那立意是不是始终在人手里？}{答“是”：收束；答“在甲方不在画师”：下一问}
\q{出于己意去创造的人，算不算创作者？}{答“算”：收束；答“不算”：点出对方把创作者窄化成接单的人}
\cut{好，我换个问法。}
\close{感谢，对方承认立意在人手里。}
\end{chain}
\begin{chain}{链二：多数链}{让对方承认好处没有只流向头部}
\q{那项研究里获赞集中度下降了，贵方否认吗？}{答“不否认”：下一问}
\q{集中度下降，是不是好处不只给了头部？}{答“是”：收束；答“那是 AI 用户”：追“AI 用户是不是人类创作者？”}
\close{感谢。好处没有只给少数人，这正是判准要看的多数。}
\end{chain}
\end{stage}

\begin{stage}{12. 正二对辩要点}{双方各 90 秒}
\textbf{任务}：守住判准，把“修图工还是立意者”这场仗打回“两半都算、看多数”。不开新战场。
\begin{points}{进攻点（4–5 个，每个 15–20 秒）}
\item 被压价的是执行还是立意？机器替得了“想画什么”吗？
\item 法院把 AI 图判给做选择的人：这是不是法律认定人的分量在选择上？
\item 获赞集中度下降：好处是不是不只给头部？
\item 贵方全场的数据量的是订单，珍视那一半贵方算过吗？
\end{points}
\begin{points}{备用点（6–8 个）}
\item 对方说用 AI 作品里的你变少 → 平均越像，有立意的人越显眼。
\item 对方说不用 AI 需要你的人变少 → 少掉的是“给我一张图”，不是“给我一个人做的图”。
\item 对方说珍视只给少数 → 标识和法律不挑名气。
\item 对方说分不出来 → 分不出来的人里八成要求标注。
\item 对方说预测到 2028 年 → 预测不是测量，而且出自利益相关方。
\item 对方说重庆工作室裁了人 → 我方承认执行岗位受冲击；题目问的是意义。
\end{points}
\begin{points}{对方最可能问的三个问题 → 回应}
\item 修图的画师立意在哪？ → 她想画什么的能力没有被拿走；今天她自己的立意不用一个团队也能做出来。
\item 珍视能养活人吗？ → 题目问的是意义，不是饭碗；我方承认处境变难。
\item 更多人创作，创作者不就更不值钱？ → 值钱的是“人的来源”相对于机器的产量；机器的产量涨得快得多。
\end{points}
\end{stage}

\begin{kit}{13. 正一结辩}{套件 · 组装后 384–408 字 / 90 秒}
\assembly{归纳分歧（任选 1） → 处理最强点（任选 1） → 回应反一（任选 1） → 临场 → 论点收束 → 价值升华 → 结尾}
\begin{pick}{归纳分歧}{1}
\begin{module}{如果全场主战场落在「需要和珍视哪个算数」}
谢谢主席。今天的主战场，落在需要和珍视哪个算数。判准说得很清楚，两半都算。谁只量需要、不把珍视放上秤，谁就只算了一半，这笔账就还没算完。
\end{module}
\begin{module}{如果全场主战场落在「修图工还是立意者」}
谢谢主席。今天的主战场，落在修图工还是立意者。机器接走的是手上的活；可想画什么、选哪一张、把哪一张交出去，从来都得由人来定。
\end{module}
\begin{fallback}{主战场不清楚时}
谢谢主席。回到判准：和 AI 爆发之前比，由人来创作这件事整体变重还是变轻。需要和珍视两半都算，而且看的是多数人。
\end{fallback}
\end{pick}
\begin{pick}{处理最强点}{1}
\begin{module}{如果反方全场最有力的是「订单和收入数据」}
反方最有力的是订单和收入。我方承认，执行型的订单在减少。可少掉的是“给我一张图”，不是“给我一个人做的图”；后一种需要，正被法律、行业和标注单独保护起来。反方量到的，只是前一种。
\end{module}
\begin{module}{如果反方全场最有力的是「修图工」}
反方最有力的，是修图工。这很疼，我方承认。但被拿走的是按星期计价的执行，不是想画什么的能力；法院判案时认的，正是人的这份选择。想画什么，机器替不了任何人。
\end{module}
\begin{fallback}{反方最有力的一点不在以上几条时}
反方最有力的一点，证明的是有人受损，这一点我方不否认。可判准问的是整体方向；受损是真的，方向却要需要和珍视两半一起算，还要看多数人，不能只看最惨的那几个人。
\end{fallback}
\end{pick}
\begin{pick}{回应反一}{1}
\begin{module}{如果反一结辩收在「珍视只落在少数人身上」}
反一刚才说，珍视只落在少数人身上。可标识约束的是机器，不挑名气；法院认的是人的选择，不是粉丝数。每一件人做的作品都算在里面。
\end{module}
\begin{module}{如果反一结辩收在「用 AI 你变少，不用 AI 需要你的人变少」}
反一刚才说，用 AI 你变少，不用 AI 需要你的人变少。还漏了一种人：用 AI 把自己想说的话说出来的人。他没有变少，他是第一次被听见。
\end{module}
\begin{module}{如果反一结辩收在一个被替代的创作者身上}
反一刚才讲了一个被替代的人。那份疼是真的。但一个人丢了活，和人来创作这件事变轻，是两回事。判准看的是多数，不是一个人。
\end{module}
\begin{fallback}{反一结辩的收束不在以上几条时}
反一刚才的结辩，也要用同一个判准去量：需要和珍视两半都算了没有，看的是不是多数人。我相信，答案各位心里已经有了。
\end{fallback}
\end{pick}
\live{40}
\begin{fixed}{论点收束}
说到底，我方两点。机器越多，人越要被认出来；执行越便宜，立意越值钱。机器接走的是手，接不走那个想说话的人。
\end{fixed}
\begin{fixed}{价值升华}
有一个人，中风以后十一年说不出完整的话。他脑子里有歌，嗓子里没有。2024 年，一首新歌用他当年的声音唱了出来，唱什么，是他和老搭档一起定的。那首歌不是机器的，是他的。我们今天争的，不是机器会不会唱歌，是每一个有话要说的人，身后是不是还站着一个“我”。
\end{fixed}
\begin{fixed}{结尾}
请各位记住：机器能做的越多，\stress{“这是一个人做的”就越重}。谢谢。
\end{fixed}
\end{kit}
\begin{contingency}
\ifthen{时间不够}{先删临场位，再把处理最强点换成兜底；价值升华和结尾不删。}
\ifthen{反一结辩收在我方没预料到的地方}{临场位只做一件事：指出它量的是一半还是两半、是少数还是多数。}
\end{contingency}

\end{document}
